% Plantilla de la casa: informe en escala de grises (estilo de la guía de Auditorita).
% Reemplaza todo lo que está entre <...>. Compila con: python <skill>/build.py <este archivo>
\documentclass[11pt,a4paper]{article}

\usepackage[spanish,es-noshorthands]{babel}
\usepackage{fontspec}
\setmainfont{texgyrepagella}[Extension=.otf, UprightFont=*-regular, BoldFont=*-bold,
  ItalicFont=*-italic, BoldItalicFont=*-bolditalic]
\setsansfont{texgyreheros}[Extension=.otf, UprightFont=*-regular, BoldFont=*-bold,
  ItalicFont=*-italic, BoldItalicFont=*-bolditalic]
\usepackage[hmargin=2.3cm, top=2.1cm, bottom=2.2cm]{geometry}
\usepackage{microtype}
\usepackage{xcolor}
\usepackage{graphicx}
\usepackage{tikz}
\usetikzlibrary{positioning,arrows.meta,fit,backgrounds,calc}
\usepackage{booktabs}
\usepackage{tabularx}
\usepackage{array}
\usepackage{needspace}
\newcolumntype{Y}{>{\raggedright\arraybackslash}X}
\usepackage{enumitem}
\usepackage{titlesec}
\usepackage{caption}
\usepackage{fancyhdr}
\usepackage{xurl}
\usepackage[hidelinks]{hyperref}

% Blanco y negro: tinta negra, acento negro y grises neutros.
\definecolor{tinta}{HTML}{111111}
\definecolor{marino}{HTML}{000000}
\definecolor{gris}{HTML}{6E6E6E}
\definecolor{linea}{HTML}{BDBDBD}
\definecolor{fondo}{HTML}{F2F2F2}
\color{tinta}

\titleformat{\section}{\Large\bfseries\color{marino}}{\thesection.}{0.6em}{}
\titleformat{\subsection}{\large\bfseries}{}{0pt}{}
\titlespacing*{\section}{0pt}{16pt}{6pt}
\setlist{itemsep=3pt,topsep=4pt,leftmargin=1.4em}
\setlength{\parindent}{0pt}
\setlength{\parskip}{6pt}
\DeclareCaptionFont{gris}{\color{gris}}
\captionsetup{font={small,it,gris}, labelfont={bf}, labelsep=period}
\renewcommand{\arraystretch}{1.25}

\pagestyle{fancy}
\fancyhf{}
\fancyfoot[C]{\small\color{gris}\thepage}
\fancyhead[R]{\small\color{gris}<Proyecto> · <qué es el documento>}
\renewcommand{\headrulewidth}{0pt}

% Diagramas: cajas grises, la caja «clave» con borde negro, pasos numerados en círculos.
\tikzset{
  caja/.style={draw=linea, fill=white, rounded corners=2pt, align=center,
    font=\sffamily\small, inner sep=6pt, minimum height=1.1cm, text=tinta},
  clave/.style={caja, draw=marino, line width=0.8pt},
  flecha/.style={-{Stealth[length=2.4mm]}, line width=0.7pt, draw=gris},
  paso/.style={circle, fill=marino, text=white, font=\sffamily\scriptsize\bfseries,
    inner sep=1.2pt, minimum size=4.2mm},
  nota/.style={font=\sffamily\footnotesize\itshape, text=gris, align=center},
}
\newcommand{\paso}[2]{\node[paso] at ($(#1.north west)+(0.05,0.02)$) {#2};}

% Recuadro gris: «En resumen» al inicio, o una decisión que el lector debe ver.
\newcommand{\recuadro}[2]{%
\noindent\fcolorbox{linea}{fondo}{\parbox{\dimexpr\linewidth-2\fboxsep-2\fboxrule}{%
\vspace{4pt}\textbf{#1}\par #2\vspace{2pt}}}\par}

\begin{document}
\thispagestyle{empty}

{\large\bfseries\color{marino} <Contexto: institución, proyecto o curso>}\par
\vspace{2pt}
{\fontsize{24}{28}\selectfont\bfseries <Título>}\par
\vspace{2pt}
{\large <Subtítulo: de qué trata, en una línea>}\par
\vspace{6pt}
{\small\color{gris} <Autor> · <día de mes de año> · <dato útil: enlace, lugar o reunión>}\par
\vspace{4pt}
{\color{linea}\rule{\linewidth}{0.8pt}}

\vspace{4pt}
\recuadro{En resumen}{%
\begin{itemize}[leftmargin=1.2em, itemsep=2pt]
  \item \textbf{<Etiqueta>.} <Una o dos frases con la cifra que importa.>
  \item \textbf{<Etiqueta>.} <...>
  \item \textbf{Qué se pide.} <La decisión o la respuesta que se espera del lector.>
\end{itemize}}

\section{<Primera sección>}

<Prosa. Cifras en formato español: 1.474 casos, 82\,\%, 0,4~h. Comillas latinas «así».>

\begin{center}
\begin{tabularx}{\linewidth}{@{}>{\bfseries}lYl@{}}
\toprule
\normalfont\bfseries <Columna> & \textbf{<Columna>} & \textbf{<Columna>} \\
\midrule
<Fila> & <Texto que puede ocupar varias líneas> & <Dato> \\
\bottomrule
\end{tabularx}
\end{center}

\begin{figure}[h]
\centering
\begin{tikzpicture}[node distance=6mm and 12mm]
  \node[caja] (a) {<Paso>\\\footnotesize\color{gris} <detalle>};
  \node[clave, right=of a] (b) {\textbf{<Paso clave>}\\\footnotesize <detalle>};
  \node[caja, right=of b] (c) {<Paso>};
  \draw[flecha] (a) -- (b);
  \draw[flecha] (b) -- (c);
  \paso{a}{1}
  \paso{b}{2}
  \paso{c}{3}
\end{tikzpicture}
\caption{<Frase que dice lo que el lector debe ver en la figura.>}
\end{figure}

\needspace{14\baselineskip}
\section{<Siguiente sección>}

<...>

\end{document}
